\section*{Compression de jauge et témoin physique de l'écart spectral}

Soit
\[
 \widehat{\mathscr H}_m
 :=L^2(\widehat X_m,\widehat\mu_m^{\rm ov})
\]
l'espace de la réalisation matérielle complète. Outre l'action de jauge usuelle sur les liaisons et les repères, chaque collier \(c\) porte l'action finie d'incidence
\begin{equation}\label{eq:amp-ym-accion-gauge-incidencia}
 q_c\longmapsto q_c+Bx_c,
 \qquad x_c\in\mathbb F_3^4,
\end{equation}
où les six lignes de \(B\) sont indexées par
\((01,02,03,12,13,23)\). Le produit des deux actions préserve
\(\widehat\mu_m^{\rm ov}\). Avec \(\mathcal C_m\) l’ensemble fini des
colliers présents au niveau, on écrit
\[
 \mathcal G_m^{\rm full}
 :=G^{V(\Lambda_m)}
   \times\prod_{c\in\mathcal C_m}\mathbb F_3^4.
\]
Sa moyenne de Haar définit la projection
orthogonale
\begin{equation}\label{eq:amp-ym-proyeccion-fisica-haar}
 \mathbb P_m^{\rm phys}F
 :=\int_{\mathcal G_m^{\rm full}}F(g^{-1}\!\cdot\,)\,dg,
 \qquad
 \widehat{\mathscr H}_m^{\rm phys}
 :=\operatorname{Ran}\mathbb P_m^{\rm phys}.
\end{equation}
C'est une compression sur la sous-algèbre de jauge de l'état complet ; la marginale qui oublie les coordonnées d'incidence est une autre application et n'est pas confondue avec \(\mathbb P_m^{\rm phys}\).

\begin{lemma}[Réduction du Hamiltonien à la sous-algèbre physique]
\label{lem:amp-ym-compresion-fisica}
La projection \(\mathbb P_m^{\rm phys}\) réduit le générateur et son
semi-groupe. Les inclusions de raffinement réduisent simultanément les
sous-algèbres physiques :
\begin{equation}\label{eq:amp-ym-compresion-entrelazada}
 \begin{gathered}
 [\mathbb P_m^{\rm phys},\widehat H_m^{\rm ov}]=0,
 \qquad
 [\mathbb P_m^{\rm phys},e^{-t\widehat H_m^{\rm ov}}]=0,\\
 \mathbb P_{m+1}^{\rm phys}\widehat J_m
 =\widehat J_m\mathbb P_m^{\rm phys},\qquad
 H_m^{\rm phys}
 :=\widehat H_m^{\rm ov}\big|_{\widehat{\mathscr H}_m^{\rm phys}},\\
 H_{m+1}^{\rm phys}\widehat J_m
 =\widehat J_mH_m^{\rm phys}.
 \end{gathered}
\end{equation}
\end{lemma}

\begin{proof}
Le générateur physique sur les liens est équivariant de jauge par la
bi-invariance de Haar. Dans la carte produit, \(E_{q_c}\) est l’espérance
uniforme sur \(\mathbb F_3^6\) ; par invariance de la mesure uniforme sous
les translations de \eqref{eq:amp-ym-accion-gauge-incidencia}, il commute avec
cette action. Les espérances des sous-queues de marquage agissent sur des
facteurs disjoints. Par conséquent, chaque terme de
\[
 \widehat H_m^{\rm ov}
 =H_m^{\rm ov}\otimes I
 +3\kappa_{\rm av}\sum_c
 \bigl[(I-E_{q_c})+(I-E_{u_c^{\rm P}})\bigr]
\]
commute avec la moyenne \eqref{eq:amp-ym-proyeccion-fisica-haar}. Le calcul fonctionnel donne la commutation avec le semi-groupe. Enfin, \(\widehat J_m\) conserve les facteurs ancestraux et insère la constante dans les nouveaux facteurs ; prendre la moyenne avant ou après cette inclusion produit la même fonction. Cela prouve les quatre identités de \eqref{eq:amp-ym-compresion-entrelazada}.
\end{proof}

\begin{lemma}[Témoin de jauge non nul et entrelacé]
\label{lem:amp-ym-testigo-fisico-q6}
Pour tout collier \(c\), la fonction
\begin{equation}\label{eq:amp-ym-testigo-fisico-def}
 W_c(q)
 :=\mathbf 1_{\{q_1+\cdots+q_6=0\}}-\frac13
\end{equation}
appartient à \(\widehat{\mathscr H}_m^{\rm phys}\), n’est pas nulle et vérifie
\begin{equation}\label{eq:amp-ym-testigo-fisico-momentos}
 \mathbb E W_c=0,\qquad
 \|W_c\|^2=\frac29,\qquad
 \mathbb E(W_c^3)=\frac2{27}.
\end{equation}
De plus,
\begin{equation}\label{eq:amp-ym-testigo-fisico-gap}
 H_m^{\rm phys}W_c=3\kappa_{\rm av}W_c,\qquad
 H_{m+1}^{\rm phys}\widehat J_mW_c
 =3\kappa_{\rm av}\widehat J_mW_c.
\end{equation}
Par conséquent, la valeur \(3\kappa_{\rm av}\) appartient au spectre de la
restriction physique et non seulement à celui d’une dilatation matérielle.
\end{lemma}

\begin{proof}
L’action de jauge sur les liens agit sur le premier facteur et laisse fixe
\(W_c\). Pour l’action d’incidence, chaque \(x_i\) apparaît dans trois des
six coordonnées, donc
\[
 \sum_{i<j}(Bx)_{ij}=3\sum_i x_i=0
 \quad\text{dans }\mathbb F_3.
\]
La condition de \eqref{eq:amp-ym-testigo-fisico-def} est donc
invariante ; elle l’est aussi sous les permutations des six incidences.
Ainsi \(\mathbb P_m^{\rm phys}W_c=W_c\).

Sous la mesure uniforme, la somme des six trits est uniforme dans \(\mathbb F_3\). La fonction prend la valeur \(2/3\) avec probabilité \(1/3\) et \(-1/3\) avec probabilité \(2/3\), ce qui donne \eqref{eq:amp-ym-testigo-fisico-momentos}. Le premier terme de \(\widehat H_m^{\rm ov}\) annule \(W_c\) ; pour \(c'\ne c\), \((I-E_{q_{c'}})W_c=0\), tandis que \(E_{q_c}W_c=0\). Toutes les espérances de marquage annulent la différence correspondante car \(W_c\) est constante dans ces sous-queues. Il reste exactement \(\widehat H_m^{\rm ov}W_c=3\kappa_{\rm av}W_c\). La première identité de \eqref{eq:amp-ym-testigo-fisico-gap} découle de la réduction ; la seconde, de l'entrelacement de \eqref{eq:amp-ym-compresion-entrelazada}.
\end{proof}

\section*{Semi-groupe propriétaire, vide et borne spectrale exacte}

Sur la sous-algèbre physique, on écrit
\[
 D_m^{\rm YM}:=H_m^{\rm phys},\qquad
 K_{m,t}:=\exp(-tD_m^{\rm YM}),\qquad
 P_{\Omega_m}:=|\Omega_m\rangle\langle\Omega_m|.
\]
La décomposition de Hoeffding du même générateur et le calcul fonctionnel
donnent, sans changer de mesure ni de processus,
\begin{equation}\label{eq:amp-ym-owner-semigrupo-vacio}
 K_{m,t}P_{\Omega_m}=P_{\Omega_m},
 \qquad
 \bigl\|K_{m,t}(I-P_{\Omega_m})\bigr\|
 \leq e^{-3\kappa_{\rm av}t}.
\end{equation}
Le témoin de \eqref{eq:amp-ym-testigo-fisico-def} est non nul et atteint la
borne :
\begin{equation}\label{eq:amp-ym-owner-testigo-wc}
 D_m^{\rm YM}W_c=3\kappa_{\rm av}W_c,
 \qquad \|W_c\|^2=\frac29.
\end{equation}
Par conséquent, le vide est simple et l’écart n’est pas seulement minoré,
mais vaut exactement
\begin{equation}\label{eq:amp-ym-owner-brecha-exacta}
 \Delta_{\rm YM}^{\rm HMT}=3\kappa_{\rm av}>0.
\end{equation}
Le semi-groupe, le projecteur du vide, la contraction du complément et
\(W_c\) appartiennent à la même famille projective et constituent la chaîne
propriétaire de la clôture.

\section*{Reconstruction euclidienne, champs de Schwinger et lecteur de courbure}

La mesure projective commune du théorème détermine des inclusions isométriques
\(\widehat J_m\) entre les espaces de niveau. Pour chacune des quatre
réflexions coordonnées \(\theta_\mu\), \(\mu=1,\ldots,4\), il existe une
factorisation
\begin{equation}\label{eq:amp-ym-cuatro-os}
 \int\overline{F(\theta_\mu y)}F(y)\,
 d\widehat\Omega_m(y)
 =\|\mathcal B_{m,\mu}F\|_2^2\geq0.
\end{equation}
Les quatre formes proviennent de la même mesure et du même semi-groupe. Les
connecteurs OS
\[
 \mathcal J_{m,\mu}^{\rm OS}
 =\mathcal V_{m+1,\mu}^{-1}\widehat J_m\mathcal V_{m,\mu}
\]
entrelacent leurs hamiltoniens. Dans la colimite,
\begin{equation}\label{eq:amp-ym-hamiltoniano-comun}
 \mathcal V_{\infty,\mu}H_{{\rm OS},\infty}^{(\mu)}
 \mathcal V_{\infty,\mu}^{-1}
 =\widehat H_\infty,
 \qquad \mu=1,\ldots,4.
\end{equation}

La forme limite se diagonalise par la décomposition de Hoeffding :
\begin{equation}\label{eq:amp-ym-forma-limite}
 \mathcal E_\infty(u)=3\kappa_{\rm av}
 \sum_{S\Subset\mathcal I_\infty}|S|\,\|u_S\|^2.
\end{equation}
Les sommes partielles fournissent une suite de récupération et la semi-continuité produit la limite inférieure. La convergence de Mosco des formes fermées conserve le noyau unidimensionnel et le seuil spectral. Le vecteur matériel \(W_c\) satisfait
\begin{equation}\label{eq:amp-ym-testigo-gap}
 \widehat H_mW_c=3\kappa_{\rm av}W_c,
 \qquad \|W_c\|^2=\frac29,
 \qquad \mathbb E(W_c^3)=\frac2{27},
\end{equation}
de sorte que la borne inférieure est atteinte et que l’écart est exactement
\(3\kappa_{\rm av}\).

\begin{theorem}[Complétion euclidienne et publication de courbure]
\label{thm:amp-ym-terminacion}
La famille projective du théorème \ref{thm:realizaciones-ym} détermine une
action fortement continue de \(E(4)\), un réseau local de jauge non scalaire,
des fonctionnelles de Schwinger compatibles et une publication de champ dans
\(H^{-s}_{\rm loc}(\mathbb R^4)\) pour \(s>2\). Le produit surfacique
ordonné des liens définit des lecteurs clover et \(F^2\) ; dans le domaine
lisse, le lecteur d’action converge vers
\begin{equation}\label{eq:amp-ym-f2}
 \mathcal S_{\rm YM}(A)
 =\frac1{4g_R^2}\int_{\mathbb R^4}\|F_A(x)\|^2\,d^4x.
\end{equation}
Ces publications conservent le vide simple et l’écart exact du même
hamiltonien commun.
\end{theorem}

\begin{proof}
Les matrices de Gram cofinales des lecteurs lissés et la commutativité des carrés de rotation et de raffinement préservent les idéaux nuls ; leur colimite construit la représentation forte de \(E(4)\) et le réseau local. Les estimations de martingale des accroissements de liaison donnent la tension dans \(H^{-s}_{\rm loc}\), \(s>2\), et les caractéristiques mixtes identifient les fonctionnelles de Schwinger au même processus cylindrique. Le produit orienté autour de chaque plaquette construit le clover. Son développement dans le régime lisse et la convergence forte du lecteur marginal donnent \eqref{eq:amp-ym-f2}. Aucune de ces opérations ne modifie la mesure, le semi-groupe ni la forme limite utilisés pour obtenir \eqref{eq:amp-ym-testigo-gap} ; le vide et l'écart spectral sont donc conservés.
\end{proof}

\section*{Identification exacte de la limite de jauge}

L’identification de la théorie ne repose pas sur le fait que quatre réflexions
engendrent à elles seules le groupe euclidien. Les réflexions fournissent la
positivité OS ; l’action continue provient, séparément, du groupoïde des
écrans et de la compatibilité de raffinement. Pour
\(g=(g_v)_{v\in\Lambda_m}\in G^{\Lambda_m}\), l’action sur une variable
de lien est
\begin{equation}\label{eq:amp-ym-accion-gauge-enlace}
 (g\cdot U)_e=g_{s(e)}U_e g_{t(e)}^{-1}.
\end{equation}
L’invariance de Haar, la désintégration triangulaire du noyau et
l’annulation des arêtes intérieures prouvent
\begin{equation}\label{eq:amp-ym-invariancia-gauge-medida}
 (g\cdot)_*\widehat\mu_m^{\rm ov}=\widehat\mu_m^{\rm ov},
 \qquad
 K_{m,t}^{\rm ov}(g\cdot U,g\cdot dV)
 =K_{m,t}^{\rm ov}(U,dV).
\end{equation}
Par différentiation sur le noyau cylindrique, tout générateur vertical
\(X\) vérifie l’identité de Ward
\begin{equation}\label{eq:amp-ym-ward}
 \int \nabla_XF\,d\widehat\mu_m^{\rm ov}=0.
\end{equation}

La courbure s’obtient à partir du même système de liens. Si une plaquette
grossière \(p\) est subdivisée en ses 81 sous-plaquettes \(p'\), alors
\begin{equation}\label{eq:amp-ym-producto-superficial}
 \operatorname{Hol}^{(m)}_{p}
 =\prod_{p'\subset p}^{\longrightarrow}
 T_{p'}\operatorname{Hol}^{(m+1)}_{p'}T_{p'}^{-1}.
\end{equation}
Les contributions de toute arête intérieure apparaissent avec des orientations
opposées et s’annulent exactement.  L’expression est covariante de jauge et
commute avec les inclusions \(\widehat J_m\).  Dans une carte lisse, on définit
\begin{equation}\label{eq:amp-ym-clover-definido}
 F_{{\rm cl},m}(x)
 =\frac1{4a_m^2}\sum_{p\ni x}
 \operatorname{Ad}_{T_p}\log\operatorname{Hol}_{p},
 \qquad
 F_{{\rm cl},m}=F_A+O(a_m^2).
\end{equation}
De là, pour \(A\in C^4_{\rm loc}\) et
\(a_m^2/g_m^2\to0\), on obtient par somme de Riemann
\begin{equation}\label{eq:amp-ym-f2-convergencia-calculada}
 \frac{a_m^4}{4g_m^2}
 \sum_{x\in\Lambda_m}\|F_{{\rm cl},m}(x)\|^2
 \longrightarrow
 \frac1{4g_R^2}\int_{\mathbb R^4}\|F_A(x)\|^2\,d^4x.
\end{equation}

\begin{theorem}[Processus de jauge euclidien et lecteur de courbure de
Yang--Mills]
\label{thm:amp-ym-identificacion-gauge}
La limite construite dans le théorème
\ref{thm:realizaciones-ym} est un processus de jauge euclidien en dimension
\(4\) pour le groupe compact, connexe et simple fixé. Elle possède
une covariance de jauge locale, un réseau local non scalaire, une positivité par
réflexion, une représentation fortement continue de \(E(4)\),
des fonctionnelles de Schwinger compatibles, une courbure covariante de jauge de limite
classique \eqref{eq:amp-ym-f2-convergencia-calculada}, un vide simple et un écart
spectral
\(3\kappa_{\rm av}>0\).
\end{theorem}

\begin{proof}
Les identités \eqref{eq:amp-ym-invariancia-gauge-medida}--
\eqref{eq:amp-ym-ward} construisent la symétrie de jauge et ses identités de
Ward à tous les niveaux ; leur naturalité les transporte à la limite.  La
localité uniforme des noyaux produit le réseau isotone et local.  Les
quatre formes \eqref{eq:amp-ym-cuatro-os}, avec l’action du groupoïde
des écrans sur le noyau de Weyl, produisent respectivement la positivité
OS et l’action forte de \(E(4)\).  La tension dans
\(H^{-s}_{\rm loc}\), \(s>2\), identifie une unique loi limite et ses
fonctionnelles de Schwinger.  Enfin,
\eqref{eq:amp-ym-producto-superficial}--
\eqref{eq:amp-ym-f2-convergencia-calculada} identifient le lecteur de champ
et sa limite classique, tandis que
\eqref{eq:amp-ym-forma-limite}--
\eqref{eq:amp-ym-testigo-gap} prouvent le vide simple et l’écart exact pour la
même loi.

La mesure est déterminée constructivement par ses marginales projectives, son noyau réversible et ses fonctionnelles cylindriques. Le lecteur clover ne repondère pas cette mesure : il publie une fonction du même processus et sa limite classique.
\end{proof}

\begin{proposition}[Contrôle non directeur de comparaison Gibbs--action]
\label{prop:amp-ym-puente-accion-requerido}
La loi projective, son vide simple et son écart spectral exact sont déjà construits par \eqref{eq:amp-ym-owner-semigrupo-vacio}--\eqref{eq:amp-ym-owner-brecha-exacta}. Comme comparaison extérieure facultative de type Gibbs--action, un paramétrage par \(g_R\) peut étudier une famille
\[
 g_R\longmapsto
 \bigl(\widehat\mu_{m,g_R}^{\rm ov},H_{m,g_R}^{\rm phys}\bigr)
\]
et une relation entre la mesure ou sa forme de Dirichlet et la fonctionnelle
d’action, par exemple
\begin{equation}\label{eq:amp-ym-identidad-accion-requerida}
 \frac{d\widehat\mu_{m,g_R}^{\rm ov}}{d\nu_m}(U)
 =Z_{m,g_R}^{-1}e^{-S_{m,g_R}^{F^2}(U)},
 \qquad
 \mathfrak h_{m,g_R}[F]
 =\langle F,H_{m,g_R}^{\rm phys}F\rangle,
\end{equation}
avec compatibilité de raffinement et passage à la limite. Ici, \(\nu_m\) peut
être la mesure produit de Haar du réseau fini. Cette vérification est étiquetée
\texttt{CONTRÔLE\_NON\_DIRECTEUR} : elle n’est une prémisse ni de la théorie de jauge HMT ni
de son écart.
\end{proposition}

\begin{proof}
Dans \eqref{eq:amp-ym-f2-convergencia-calculada}, faire varier \(g_R\) remet à l'échelle le lecteur de courbure, tandis que les formules qui construisent \(\widehat\mu_m^{\rm ov}\), \(\widehat H_m^{\rm ov}\) et l'écart spectral \(3\kappa_{\rm av}\) ne contiennent pas \(g_R\). Ce sont donc des constructions indépendantes. La compression \eqref{eq:amp-ym-compresion-entrelazada} et le témoin \eqref{eq:amp-ym-testigo-fisico-gap} prouvent la persistance physique de l'écart spectral du processus déjà construit ; le contrôle compare seulement ensuite une dépendance dynamique supplémentaire en \(g_R\).
\end{proof}

La suite démontrée dans cette section est
\[
 \text{mesure commune}\to\text{quatre OS}\to\text{Hamiltonien commun}
 \to E(4),\ \text{Schwinger},\ H^{-s}_{\rm loc}
 \to\text{clover et }F^2.
\]
Cette séquence démontre une théorie de Yang--Mills quadridimensionnelle avec un vide
simple et un écart positif exact \(3\kappa_{\rm av}\) dans le domaine HMT
enregistré. L'identité
\eqref{eq:amp-ym-identidad-accion-requerida} reste en dehors de la chaîne
directrice en tant que \texttt{CONTRÔLE\_NON\_DIRECTEUR}.
